% summary-preamble.tex
% Included in the plain-language summary build via --include-in-header.
% Intentionally omits \sectionbreak (\clearpage) — the summary is a
% shorter, reader-friendly document where forced page-breaks per section
% would be unnecessarily disruptive.

% Unicode character mappings for XeLaTeX
% The plain-language summary uses Unicode symbols directly in text mode:
%   - Proof-step marker: ✓
%   - Comparison operators: ≠ ≡ ≈ ≥ ≤
%   - Greek letter shorthand: δ
%   - Unicode superscripts for exponents: ⁰¹²³⁴⁵⁶⁷⁸⁹⁻
%     (e.g. 10⁻⁷, 10²⁹, 10³⁵ rather than LaTeX math mode)
% lmroman lacks these glyphs; newunicodechar routes them to proper LaTeX
% commands so no glyph-missing warnings are emitted.
\usepackage{amssymb}
\usepackage{newunicodechar}

% Typography: paragraph spacing and line height for readability.
% parskip replaces first-line indents with inter-paragraph gaps.
% setspace opens up line height slightly (1.2x, not 1.5x).
% \usepackage{parskip}
% \usepackage{setspace}
% \setstretch{1.2}

% Section spacing: generous gap above headings to visually separate
% sections; tighter gap below so the heading stays close to its body.
%   \titlespacing*{command}{left}{before-sep}{after-sep}
%   body → \section:      4 lines  (clear section break)
%   body → \subsection:   2.5 lines (clear but smaller)
%   \section → \subsection: 0.5 + 2.5 = 3 lines (compact heading stack)
%   heading → body text:  0.5 lines (heading stays with its content)
\usepackage{titlesec}
\titlespacing*{\section}{0pt}{4\baselineskip}{0.5\baselineskip}
\titlespacing*{\subsection}{0pt}{2.5\baselineskip}{0.5\baselineskip}
\titlespacing*{\subsubsection}{0pt}{2\baselineskip}{0.5\baselineskip}
% Proof marker and comparison operators
\newunicodechar{✓}{$\checkmark$}
\newunicodechar{≠}{$\neq$}
\newunicodechar{≡}{$\equiv$}
\newunicodechar{≈}{$\approx$}
\newunicodechar{≥}{$\geq$}
\newunicodechar{≤}{$\leq$}
% Greek shorthand
\newunicodechar{δ}{$\delta$}
% Unicode superscript digits and minus (used for exponents like 10⁻⁷, 10²⁹)
\newunicodechar{⁰}{$^{0}$}
\newunicodechar{¹}{$^{1}$}
\newunicodechar{²}{$^{2}$}
\newunicodechar{³}{$^{3}$}
\newunicodechar{⁴}{$^{4}$}
\newunicodechar{⁵}{$^{5}$}
\newunicodechar{⁶}{$^{6}$}
\newunicodechar{⁷}{$^{7}$}
\newunicodechar{⁸}{$^{8}$}
\newunicodechar{⁹}{$^{9}$}
\newunicodechar{⁻}{$^{-}$}
